% WRONG This thesis didnt validate to what extent can gradient projection methods be
% implemented for hypernetworks. THAT WAS THE FIRST PREMISE. But I failed to create a
% principled normalization structure for this combination. 

% NOW the rq1 looks like this:
% \item[\textbf{Sub-RQ1}] \textbf{(Synergy vs.\ Conflict):} Does the joint
% integration of hard gradient projection with the hypernetwork's soft
% functional regularizer yield a constructive
% synergy, or does a standalone regularized hypernetwork offer a more
% efficient optimization frontier?

% So basically I didnt resolve that question. I made one attempt but it basically failed
% to improve the performance of hypernetworks as plain adam with hypernetwork outperformed every 
% other method

\section{Conclusion}
\label{sec:conclusion}

This thesis investigated integrating gradient projection
methods into chunked hypernetworks for continual learning.
We primarily diagnosed the architectural boundaries
and numerical pathologies governing this core integration.
Alongside these fundamental structural findings, we
also evaluated adjacent methodological optimizations regarding
parameter inertia and Fisher matrix accumulation.

\paragraph{Sub-RQ1 --- Partial Synergy and the Momentum Bottleneck.}
The synergy between hard gradient projection and
soft functional regularization is constructive but bounded.
While \texttt{iFOPNG} outperforms unconstrained \texttt{SGD}
in standard-capacity settings, standalone regularized
hypernetworks optimized via \texttt{Adam} remain
vastly superior.

This performance gap stems from an architectural
conflict, not projection geometry.
Under extreme compression ($d_h = 8$), Fisher-based
projection methods suffer a severe \emph{retention}
failure, experiencing substantial negative backward transfer
as the constraint fails to protect prior knowledge.
Because projection modifies raw gradients before the
optimizer updates, it discards adaptive momentum.
In hypernetworks, this momentum acts as a crucial
structural regularizer; without it, the network commits
prematurely to poorly-conditioned gradient directions,
destabilizing the generative manifold and causing
catastrophic interference.

Ultimately, the loss of adaptive momentum is
the binding constraint.
Integrating gradient projection into hypernetworks
will require projection-compatible adaptive optimizers,
rather than post-hoc gradient modifications \citep{hu_hidden_2026}.

\paragraph{Sub-RQ2 --- Normalization Is a Necessary Condition.}
Without explicit local rescaling, projection methods are numerically
incompatible with chunked hypernetworks: matrix inversion in the
projection kernel degrades or collapses entirely depending on chunk
count, pulling performance below plain SGD.
QR orthonormalization of the gradient basis combined with
absolute-maximum Fisher scaling resolves the pathology in the tested
configurations.
This normalization scheme is a necessary engineering precondition for
any future combination of gradient projection with chunked
meta-learners, regardless of which projection algebra is employed.

\paragraph{Sub-RQ3 --- Parameter Inertia via Combined Fisher Metric.}
Replacing the task-specific Fisher metric with the combined inertial metric
($F_c$) improves retention in direction across every benchmark tested.
The advantage is statistically significant on four of the five standalone
benchmarks, including Permuted-MNIST, where it is negligible in size
($+0.3$~pp) yet reliable on account of its consistency across seeds; the only
standalone null is the near-ceiling Split-MNIST multi-head ($+0.7$~pp,
$p = .325$). The advantage is largest on Split-CIFAR10 ($+14.1$~pp,
$p = .024$), alongside a substantial gain on Split-MNIST single-head ($+1.3$~pp). 
The same positive direction carries into the three hypernetwork
settings, though none reaches significance under the elevated inter-seed
variance of that regime. Integrating parameter inertia therefore yields a
positive retention advantage in every setting: statistically robust across the
standalone benchmarks, and directionally consistent but underpowered in the
hypernetwork settings.

\paragraph{Sub-RQ4 --- MAX Accumulation Is the Clear Practical Choice.}
On a twenty-task sequence the element-wise maximum operator was statistically
superior to an exponential moving average (${\sim}1.60$~pp, $p < 0.001$),
and the hypothesised plasticity cost of its monotonic growth did not
materialise. MAX is therefore preferred: hyperparameter-free, simpler, and
marginally better.

\paragraph{Implications for the field.}
The results establish a clear empirical boundary for the current generation
of gradient projection methods: they can be made numerically stable in
chunked hypernetwork settings via local rescaling, and their Fisher geometry
can be improved via the combined Fisher metric $F_c$, but the momentum incompatibility
prevents them from reaching the performance frontier of adaptive
unconstrained optimizers in generative parameter spaces.
The normalization contribution is directly applicable to any method that
maintains a gradient subspace memory alongside a chunked weight generator.
The iFOPNG contribution extends naturally to any continual learning
setting where Fisher overlap accumulates across tasks.
Together, they narrow the gap between projection-based and
regularization-based approaches, and provide a principled foundation for
the next step: a projected adaptive optimizer that respects both the
geometric constraints of the projection and the curvature structure that
adaptive moment estimates capture.
